\notepage{55}{Ehrenfest's theorem and the virial theorem}
For an operator $\Omega$ with no explicit time dependence, the
Heisenberg equation implies
\[
 \ii\hbar\frac{\dd}{\dd t}\langle\Omega\rangle
 =\langle\comm{\Omega}{H}\rangle.
\]
In an energy eigenstate,
\begin{align*}
 \mel{E}{\comm{\Omega}{H}}{E}
 &=\mel{E}{\Omega H}{E}-\mel{E}{H\Omega}{E}\\
 &=E\bigl(\mel{E}{\Omega}{E}-\mel{E}{\Omega}{E}\bigr)=0.
\end{align*}
Consequently all well-defined expectation values of such operators are
stationary in an energy eigenstate.

\smallheading{Choose \texorpdfstring{$\Omega=\mathbf R\cdot\mathbf P$}{Omega=R dot P}}
For $H=\mathbf P^2/(2m)+V(\mathbf R)$, the operator equations give
\begin{align*}
 \dot X_i&=\frac1{\ii\hbar}\comm{X_i}{H}
   =\frac1{2m\ii\hbar}\comm{X_i}{P_i^2}=\frac{P_i}{m},\\
 \dot P_i&=\frac1{\ii\hbar}\comm{P_i}{H}
   =\frac1{\ii\hbar}\comm{P_i}{V}
   =-\frac{\partial V}{\partial X_i}.
\end{align*}
In vector notation, $\dot{\mathbf R}=\mathbf P/m$ and
$\dot{\mathbf P}=-\boldsymbol\nabla_RV$. Therefore a stationary state
satisfies
\[
 0=\left\langle\dot{\mathbf R}\cdot\mathbf P\right\rangle
   +\left\langle\mathbf R\cdot\dot{\mathbf P}\right\rangle
 =\frac{\langle\mathbf P^2\rangle}{m}
  -\langle\mathbf R\cdot\boldsymbol\nabla_RV\rangle.
\]
If $V(\mathbf R)=cR^k=c(X_1^2+X_2^2+X_3^2)^{k/2}$, then
\[
 -\frac{\partial V}{\partial X_i}=-ckX_iR^{k-2},\qquad
 \mathbf R\cdot\boldsymbol\nabla_RV=ckR^k=kV.
\]
Hence
\[
 \boxed{\langle T\rangle
 =\left\langle\frac{\mathbf P^2}{2m}\right\rangle
 =\frac{k}{2}\langle V\rangle.}
\]
This is the virial theorem for a homogeneous power-law potential,
for stationary states where the indicated expectations and boundary
conditions are well defined.
